% GENERATED FILE. Edit canonical lessons, then run npm run generate:books.
% canonical-source-hash: fnv1a64:da9581237fe54009
% canonical-lessons: ES-C678-visitar, ES-C678-ermita, ES-C678-estatua, ES-C678-torre, ES-C678-galeria

\chapter{De visita}
\label{ch:es-de-visita}
\hlchaptermodality{678}

\begin{chapteropening}
\textbf{By the end of this chapter:} \emph{I can plan a visit to a town: statues, towers, galleries and a country chapel.}

Complete the last lesson of chapter 678: I can plan a visit to a town: statues, towers, galleries and a country chapel.
\end{chapteropening}

\section[visitar]{visitar — to visit}

\label{lesson:ES-C678-visitar}



\begin{quote}
Before the new one: say the Spanish for the palace, then the Spanish for the city wall.
\end{quote}

\subsection*{You'll want to know: visitar}
\textbf{visitar} — \textquotedblleft{}to visit\textquotedblright{}.

\begin{quote}
Queremos visitar el museo mañana.
\end{quote}

\begin{quote}
\emph{We want to visit the museum tomorrow.}
\end{quote}

For people, use \emph{a}: \emph{visitar a mis abuelos}.

\begin{cousinweb}
\emph{visitar} comes from Latin \emph{vīsitāre} \textquotedblleft{}to see, to go to see\textquotedblright{} (Lewis \& Short; EtymDB), built on \emph{vīsere} \textquotedblleft{}to go and see\textquotedblright{}, from \emph{vidēre} \textquotedblleft{}to see\textquotedblright{}.
\end{cousinweb}

\begin{grammarlens}[title={What you've built}]
One more word for telling, explaining and describing at B1.
\end{grammarlens}

\subsection*{Your turn}
\begin{itemize}
\raggedright
  \item \emph{Say it:} \emph{visitar}
  \item \emph{Say it:} \emph{visitar}, once more
  \item \emph{Say it:} a sentence of your own with \emph{visitar}
  \item \emph{Recall:} say the Spanish for the palace, then the Spanish for the city wall, then say \emph{visitar} again
\end{itemize}

\begin{tcolorbox}[breakable,colback=teal!4,colframe=teal!35!black,title={Before you move on}]
What does visitar mean? (\textquotedblleft{}To visit\textquotedblright{}.) Say it once more.
\end{tcolorbox}
\section[la ermita]{la ermita — the country chapel, the hermitage}

\label{lesson:ES-C678-ermita}



\begin{quote}
Before the new one: say the Spanish for the city wall, then the Spanish for to visit.
\end{quote}

\subsection*{You'll want to know: la ermita}
\textbf{la ermita} — \textquotedblleft{}the country chapel, the hermitage\textquotedblright{}.

\begin{quote}
Subimos a la ermita del pueblo.
\end{quote}

\begin{quote}
\emph{We walked up to the village chapel.}
\end{quote}

Many villages walk to their \emph{ermita} on a feast day.

\begin{cousinweb}
\emph{ermita} goes back to Late Latin \emph{erēmīta} \textquotedblleft{}a hermit\textquotedblright{} (Lewis \& Short), from Greek \emph{erēmítēs}, built on \emph{érēmos} \textquotedblleft{}lonely, desert\textquotedblright{}: the hermit lived alone.
\end{cousinweb}

\begin{grammarlens}[title={What you've built}]
One more word for telling, explaining and describing at B1.
\end{grammarlens}

\subsection*{Your turn}
\begin{itemize}
\raggedright
  \item \emph{Say it:} \emph{la ermita}
  \item \emph{Say it:} \emph{la ermita}, once more
  \item \emph{Say it:} a sentence of your own with \emph{la ermita}
  \item \emph{Recall:} say the Spanish for the city wall, then the Spanish for to visit, then say \emph{la ermita} again
\end{itemize}

\begin{tcolorbox}[breakable,colback=teal!4,colframe=teal!35!black,title={Before you move on}]
What does la ermita mean? (\textquotedblleft{}The country chapel, the hermitage\textquotedblright{}.) Say it once more.
\end{tcolorbox}
\section[la estatua]{la estatua — the statue}

\label{lesson:ES-C678-estatua}



\begin{quote}
Before the new one: say the Spanish for to visit, then the Spanish for the country chapel.
\end{quote}

\subsection*{You'll want to know: la estatua}
\textbf{la estatua} — \textquotedblleft{}the statue\textquotedblright{}.

\begin{quote}
Hay una estatua en el centro de la plaza.
\end{quote}

\begin{quote}
\emph{There is a statue in the middle of the square.}
\end{quote}

\emph{Quedarse como una estatua} — to stand stock-still.

\begin{cousinweb}
\emph{estatua} is a learned borrowing from Latin \emph{statua} \textquotedblleft{}an image, a statue\textquotedblright{} (Lewis \& Short; EtymDB), from \emph{statuere} \textquotedblleft{}to set up\textquotedblright{}.
\end{cousinweb}

\begin{grammarlens}[title={What you've built}]
One more word for telling, explaining and describing at B1.
\end{grammarlens}

\subsection*{Your turn}
\begin{itemize}
\raggedright
  \item \emph{Say it:} \emph{la estatua}
  \item \emph{Say it:} \emph{la estatua}, once more
  \item \emph{Say it:} a sentence of your own with \emph{la estatua}
  \item \emph{Recall:} say the Spanish for to visit, then the Spanish for the country chapel, then say \emph{la estatua} again
\end{itemize}

\begin{tcolorbox}[breakable,colback=teal!4,colframe=teal!35!black,title={Before you move on}]
What does la estatua mean? (\textquotedblleft{}The statue\textquotedblright{}.) Say it once more.
\end{tcolorbox}
\section[la torre]{la torre — the tower}

\label{lesson:ES-C678-torre}



\begin{quote}
Before the new one: say the Spanish for the country chapel, then the Spanish for the statue.
\end{quote}

\subsection*{You'll want to know: la torre}
\textbf{la torre} — \textquotedblleft{}the tower\textquotedblright{}.

\begin{quote}
Subimos a la torre de la iglesia.
\end{quote}

\begin{quote}
\emph{We went up the church tower.}
\end{quote}

In Spain \emph{una torre} is also a tall block of flats.

\begin{cousinweb}
\emph{torre} is inherited, through Old Spanish \emph{torre}, from Latin \emph{turris} \textquotedblleft{}a tower\textquotedblright{} (Lewis \& Short; EtymDB).
\end{cousinweb}

\begin{grammarlens}[title={What you've built}]
One more word for telling, explaining and describing at B1.
\end{grammarlens}

\subsection*{Your turn}
\begin{itemize}
\raggedright
  \item \emph{Say it:} \emph{la torre}
  \item \emph{Say it:} \emph{la torre}, once more
  \item \emph{Say it:} a sentence of your own with \emph{la torre}
  \item \emph{Recall:} say the Spanish for the country chapel, then the Spanish for the statue, then say \emph{la torre} again
\end{itemize}

\begin{tcolorbox}[breakable,colback=teal!4,colframe=teal!35!black,title={Before you move on}]
What does la torre mean? (\textquotedblleft{}The tower\textquotedblright{}.) Say it once more.
\end{tcolorbox}
\section[la galería]{la galería — the gallery; the covered balcony}

\label{lesson:ES-C678-galeria}



\begin{quote}
Before the new one: say the Spanish for the statue, then the Spanish for the tower.
\end{quote}

\subsection*{You'll want to know: la galería}
\textbf{la galería} — \textquotedblleft{}the gallery; the covered balcony\textquotedblright{}.

\begin{quote}
La galería tiene cuadros modernos.
\end{quote}

\begin{quote}
\emph{The gallery has modern paintings.}
\end{quote}

\emph{Galerías comerciales} — a shopping arcade.

\begin{cousinweb}
\emph{galería} came from Italian \emph{galleria}, a covered walk; where the word itself comes from is not settled.
\end{cousinweb}

\begin{grammarlens}[title={What you've built}]
One more word for telling, explaining and describing at B1.
\end{grammarlens}

\subsection*{Your turn}
\begin{itemize}
\raggedright
  \item \emph{Say it:} \emph{la galería}
  \item \emph{Say it:} \emph{la galería}, once more
  \item \emph{Say it:} a sentence of your own with \emph{la galería}
  \item \emph{Recall:} say the Spanish for the statue, then the Spanish for the tower, then say \emph{la galería} again
\end{itemize}

\begin{tcolorbox}[breakable,colback=teal!4,colframe=teal!35!black,title={Before you move on}]
What does la galería mean? (\textquotedblleft{}The gallery; the covered balcony\textquotedblright{}.) Say it once more.
\end{tcolorbox}
